

\begin{abstract}[it]
Este proyecto estima la pobreza multidimensional a nivel municipal en Colombia mediante modelos de Estimación en Áreas Pequeñas (SAE) complementados con enfoques espaciales y fuentes oficiales del DANE, en concordancia con el principio de «no dejar a nadie atrás» de la Agenda 2030. La importancia del estudio radica en que el país carece de estimaciones intercensales continuas del Índice de Pobreza Multidimensional (IPM) para los 1.122 municipios, debido a las restricciones de representatividad muestral que sufren las encuestas de hogares tradicionales al desagregarse a escala subnacional. Para mitigar esta limitación, el objetivo general del proyecto consistió en desarrollar y comparar formulaciones de modelos SAE, incorporando y evaluando componentes espaciales para estimar el IPM municipal a partir de fuentes oficiales. La metodología abordó la extensión del modelo de área de Fay-Herriot mediante transformaciones no lineales acotadas que permitan obtener resultados dentro del intervalo [0,1], además de la integración de una propuesta de estructura de dependencia espacial Condicional Autorregresiva (CAR), articulando los microdatos encuestales con registros administrativos e indicadores geográficos auxiliares. Como resultado de la evaluación de desempeño, se demostró que la incorporación del componente espacial CAR en el marco Fay-Herriot puede optimizar la precisión de las estimaciones en municipios con baja representatividad muestral y capturar con éxito la autocorrelación geográfica intermunicipal. Asimismo, el modelado permitió identificar las variables auxiliares oficiales del DANE con mayor poder explicativo sobre las privaciones multidimensionales. Finalmente, las estimaciones indirectas generadas permitieron caracterizar la distribución territorial y las estructuras de dependencia espacial del IPM municipal, visibilizando patrones de aglomeración y desigualdades territoriales. Este trabajo aporta una metodología estadística reproducible que podría servir de base a entidades nacionales y locales para fortalecer el diagnóstico socioeconómico intercensal y optimizar la focalización de la inversión pública.
\end{abstract}